% !TEX root = ../main.tex
\section{Aplicación de la orientación a objetos y asignación de responsabilidades}\label{sec:responsabilidades}

Esta sección justifica el diseño de clases desde dos perspectivas complementarias: los principios de la programación orientada a objetos y los principios de asignación de responsabilidades propuestos por \citet{larman2003uml}. Para cada decisión relevante se indica la alternativa descartada.

\subsection{Principios de la orientación a objetos}

\begin{table}[H]
\caption{Aplicación de los principios de la orientación a objetos}\label{tab:poo}
\tablacompacta
\begin{tabularx}{\textwidth}{@{}L{2.9cm}Y@{}}
  \toprule
  \cab{Principio} & \cab{Aplicación en el diseño} \\ \midrule
  Abstracción & \texttt{Jugador}, \texttt{EstadoJuego} y \texttt{Jugada} son abstractas e \texttt{IEstrategiaCpu} es una interfaz: definen lo que ofrece cada familia sin fijar cómo lo hace. \texttt{IReloj} e \texttt{IRepositorioPartidas} abstraen la hora y el almacenamiento. \\ \filasep
  Encapsulamiento & \texttt{Mano}, \texttt{Mesa}, \texttt{FichasParaRobar} y \texttt{HistorialJugadas} ocultan sus colecciones internas y solo permiten modificarlas mediante operaciones que preservan las reglas. Por ejemplo, ninguna clase externa puede agregar una ficha a la mesa sin pasar por \texttt{PuedeColocar}. \texttt{Ficha} es inmutable. \\ \filasep
  Herencia & \texttt{JugadorHumano} y \texttt{JugadorCpu} heredan de \texttt{Jugador}; los cuatro estados heredan de \texttt{EstadoJuego} y las cuatro jugadas de \texttt{Jugada}. Los niveles de CPU dejaron de ser subclases: se componen como estrategias (sección~\ref{sec:patrones}). \\ \filasep
  Polimorfismo & \texttt{Partida} consulta \texttt{AplicaLimiteTiempo} sin distinguir el tipo concreto de jugador, delega en su estado sin conocer cuál es y ejecuta o deshace jugadas sin conocer su tipo. \texttt{JugadorCpu} invoca \texttt{Elegir} sin conocer el nivel de su estrategia. \\
  \bottomrule
\end{tabularx}
\end{table}

\clearpage
\subsection{Principios de asignación de responsabilidades}

{\tablacompacta
\begin{longtable}{@{}L{2.7cm}L{4.3cm}L{7.5cm}@{}}
  \caption{Justificación de las responsabilidades según los principios GRASP}\label{tab:grasp}\\
  \toprule
  \cab{Principio} & \cab{Decisión} & \cab{Justificación y alternativa descartada} \\ \midrule \endfirsthead
  \toprule
  \cab{Principio} & \cab{Decisión} & \cab{Justificación y alternativa descartada} \\ \midrule \endhead
  \bottomrule \endlastfoot
  Controlador & \texttt{ControladorJuego} recibe todas las operaciones del sistema. & Es el punto único de entrada desde la interfaz, de modo que la consola no conoce \texttt{Partida} por dentro. Se descartó que la consola invocara directamente a \texttt{Partida}, porque acoplaría la interfaz a las reglas y obligaría a repetir la coordinación en la simulación. \\ \filasep
  Experto en información & \texttt{Ficha} calcula sus puntos y si es doble; \texttt{Mano} suma sus puntos y encuentra su mayor doble; \texttt{Mesa} decide si una ficha puede colocarse. & Cada responsabilidad queda en la clase que posee los datos necesarios. Se descartó concentrar estos cálculos en \texttt{Partida}, que tendría que recorrer las colecciones de otras clases y conocer su estructura interna. \\ \filasep
  Experto en información & \texttt{Resultado} determina el ganador y el desempate al cerrarse el juego (RN-09). & El resultado conoce los puntajes revelados de todos los jugadores. Se descartó ubicar el desempate en \texttt{ControladorJuego}, que no posee esa información. \\ \filasep
  Creador & \texttt{Partida} crea \texttt{Mesa}, \texttt{HistorialJugadas} y los jugadores. & La partida los contiene por composición y dispone de los datos de la configuración para inicializarlos. \\ \filasep
  Creador y fabricación pura & \texttt{Repartidor} crea las fichas y las fichas para robar. & Barajar y repartir no es responsabilidad natural de ningún concepto del juego. Una clase dedicada conserva la cohesión de \texttt{Partida} y permite fijar la semilla del barajado en la simulación. \\ \filasep
  Fabricación pura & \texttt{HistorialJugadas} administra la secuencia de jugadas para deshacer. & Separa de \texttt{Partida} la lógica de localizar la última jugada del humano y las jugadas de CPU posteriores (RN-10). \\ \filasep
  Fabricación pura & \texttt{VistaPartida} expone solo la información visible. & Controla en un único lugar qué ve cada jugador, incluida la revelación de fichas al cerrarse el juego. Solo contiene tipos de valor, de modo que la interfaz no puede modificar el dominio a través de ella. Se descartó entregar \texttt{Partida} a la interfaz y a las CPU, porque expondría las manos ajenas. \\ \filasep
  Fabricación pura & \texttt{PartidaGuardada} contiene la configuración, la semilla y las jugadas. & Permite guardar la partida sin exponer el estado interno de \texttt{Partida}: al cargarla, la partida se reconstruye repitiendo el reparto y las jugadas. Se descartó serializar el objeto \texttt{Partida} completo, porque obligaría a hacer públicos sus campos y acoplaría el formato del archivo a su estructura interna. \\ \filasep
  Polimorfismo & Jerarquías de jugadores, estrategias, estados y jugadas. & El comportamiento que varía según el tipo de jugador, el nivel, el estado de la partida o el tipo de jugada se resuelve con operaciones polimórficas. Se descartaron atributos de tipo con condicionales en \texttt{Partida}, que crecerían con cada nuevo caso. Los patrones de la sección~\ref{sec:patrones} son la forma concreta de este principio. \\ \filasep
  Variaciones protegidas e indirección & \texttt{IReloj} e \texttt{IRepositorioPartidas}. & Protegen al dominio de los cambios en la fuente de la hora y en el medio de almacenamiento. Permiten que la simulación use un reloj controlado. Se descartó usar directamente \texttt{DateTime.Now} y el sistema de archivos desde el dominio. \\ \filasep
  Bajo acoplamiento & La consola depende solo de \texttt{ControladorJuego} y \texttt{VistaPartida}; el dominio no depende de ningún otro proyecto. & Cumple la restricción de independencia del dominio y permite reemplazar la interfaz sin modificar las reglas. \\ \filasep
  Alta cohesión & Cada clase concentra un único tipo de responsabilidad. & \texttt{Partida} coordina las reglas, mientras que el reparto, el historial, la vista y el cálculo de resultados residen en clases propias. \\
\end{longtable}}

\subsection{Puntos de variación identificados}

La versión 1.0 dejó identificados cuatro puntos donde el juego varía: la forma en que cada nivel de CPU decide su jugada, el comportamiento de la partida según su estado, la reversión de jugadas y la construcción de la partida a partir de la configuración. Los tres primeros se resolvieron en esta entrega con Strategy, State y Command (sección~\ref{sec:patrones}). La construcción de la partida y de sus jugadores corresponde a los patrones creacionales de la Entrega~3.
